\subsection{Implementazione del modello}

Il modello è implementato mediante PyTorch e organizzato in tre componenti principali: encoder convoluzionale, nucleo ricorsivo e Policy Head. Queste componenti sono definite in moduli separati e integrate dalla classe \texttt{GoTRMModel}, implementata in \texttt{go\_model.py}.

\subsubsection{Encoder convoluzionale}

L'encoder è implementato nella classe \texttt{GoEncoder}, definita in \texttt{go\_encoder.py}. Il blocco convoluzionale iniziale è costituito da una convoluzione bidimensionale, seguita da Batch Normalization e ReLU:

\begin{lstlisting}[language=Python, caption={Definizione del blocco convoluzionale di GoEncoder}, label={lst:go_encoder}]
self.conv_stem = nn.Sequential(
    nn.Conv2d(in_channels, hidden_size,
              kernel_size=3, padding=1, bias=False),
    nn.BatchNorm2d(hidden_size),
    nn.ReLU()
)
\end{lstlisting}

La convoluzione trasforma i sette canali iniziali in 256 mappe di caratteristiche, mantenendo invariate le dimensioni spaziali grazie al kernel \(3\times3\) e al padding pari a 1.

Il risultato viene successivamente convertito nella sequenza di token richiesta dal nucleo ricorsivo. Le dimensioni spaziali vengono appiattite e riordinate, mentre a ciascun token viene sommato il corrispondente embedding posizionale apprendibile:

\begin{lstlisting}[language=Python, caption={Conversione delle mappe convoluzionali in una sequenza}, label={lst:go_encoder_sequence}]
x = self.conv_stem(x)

batch_size = x.size(0)
x = x.view(batch_size, x.size(1), -1)
x = x.transpose(1, 2)

x = x + self.pos_embedding
\end{lstlisting}

L'output assume forma \([B,N^2,256]\), dove \(B\) indica la dimensione del batch e \(N\) quella della scacchiera.

\subsubsection{Nucleo ricorsivo}

Il nucleo è implementato nel modulo \texttt{go\_trm.py}, che contiene le classi dedicate alla configurazione, ai blocchi di elaborazione e alla gestione degli stati ricorsivi.

Gli stati \(z_H\) e \(z_L\) sono rappresentati mediante tensori di forma \([B,N^2,256]\). La loro inizializzazione e il ripristino sono gestiti rispettivamente dai metodi \texttt{empty\_carry()} e \texttt{reset\_carry()}.

Il metodo \texttt{forward()} coordina gli aggiornamenti ricorsivi, reinserendo la rappresentazione prodotta dall'encoder durante gli aggiornamenti di \(z_L\):

\begin{lstlisting}[language=Python, caption={Aggiornamento ricorsivo degli stati latenti}, label={lst:trm_cycles}]
input_embeddings = self.go_encoder(batch["inputs"])

z_H, z_L = carry.z_H, carry.z_L

with torch.no_grad():
    for _ in range(self.config.H_cycles - 1):
        for _ in range(self.config.L_cycles):
            z_L = self.L_level(z_L, z_H + input_embeddings)
        z_H = self.L_level(z_H, z_L)

for _ in range(self.config.L_cycles):
    z_L = self.L_level(z_L, z_H + input_embeddings)
z_H = self.L_level(z_H, z_L)
\end{lstlisting}

Gli aggiornamenti utilizzano lo stesso modulo \texttt{L\_level}, condividendone i parametri tra i due stati e le diverse iterazioni.

La gestione del gradiente costituisce un aspetto rilevante dell'implementazione: i primi \(H\_cycles-1\) cicli vengono eseguiti mediante \texttt{torch.no\_grad()}, mentre il ciclo finale mantiene attivo il calcolo necessario alla retropropagazione. Al termine, gli stati aggiornati vengono salvati nella struttura \texttt{carry}, utilizzando \texttt{detach()} per separarli dal grafo computazionale.

Il numero di aggiornamenti è determinato dai parametri \texttt{H\_cycles} e \texttt{L\_cycles}, consentendo di modificare la profondità ricorsiva attraverso la configurazione senza intervenire sul codice del modello.

\subsubsection{Policy Head}

La classe \texttt{GoPolicyHead}, definita in \texttt{go\_policy\_head.py}, riceve la sequenza prodotta dal nucleo ricorsivo e genera i logits associati alle azioni.

Per le intersezioni viene utilizzato un MLP condiviso tra tutti i token:

\begin{lstlisting}[language=Python, caption={MLP per la previsione delle mosse sul tabellone}, label={lst:policy_board}]
self.board_mlp = nn.Sequential(
    nn.Linear(hidden_size, hidden_size // 2),
    nn.ReLU(),
    nn.Linear(hidden_size // 2, 1)
)
\end{lstlisting}

Il ramo dedicato a \textit{PASS} utilizza invece una rappresentazione globale, ottenuta mediante media delle rappresentazioni latenti lungo la dimensione della sequenza:

\begin{lstlisting}[language=Python, caption={Generazione del logit PASS}, label={lst:policy_pass}]
global_features = x.mean(dim=1)
pass_logit = self.pass_mlp(global_features)
\end{lstlisting}

I risultati vengono infine concatenati per produrre un vettore di \(N^2+1\) valori:

\begin{lstlisting}[language=Python, caption={Composizione dell'output della Policy Head}, label={lst:policy_output}]
board_logits = self.board_mlp(x).squeeze(-1)
logits = torch.cat([board_logits, pass_logit], dim=-1)
return logits
\end{lstlisting}

I primi \(N^2\) valori corrispondono alle intersezioni, mentre l'ultimo rappresenta l'azione \textit{PASS}. I valori ottenuti sono logits non normalizzati e vengono convertiti in probabilità soltanto durante la fase di inferenza.

\paragraph{Integrazione dei componenti}

La classe \texttt{GoTRMModel}, definita in \texttt{go\_model.py}, integra encoder, nucleo ricorsivo e Policy Head. Il metodo \texttt{forward()} coordina il passaggio delle rappresentazioni tra i moduli:

\begin{lstlisting}[language=Python, caption={Integrazione tra nucleo ricorsivo e Policy Head}, label={lst:go_model}]
carry = self.trm.initial_carry(batch_dict)
carry, hidden_features = self.trm(carry, batch_dict)
return self.policy_head(hidden_features)
\end{lstlisting}

In questo modo, le procedure di addestramento e inferenza possono utilizzare il modello completo attraverso un'unica interfaccia, senza dover gestire separatamente le singole componenti.
